%======================================================================
% Chapter 1: Numbers, Arithmetic, and Algebraic Thinking
%======================================================================
\chapter{Numbers, Arithmetic, and Algebraic Thinking}\label{ch:numbers}

\begin{motivation}
To a computer, everything in machine learning is a number. A photograph of a
cat is a grid of numbers that record how bright each tiny dot (pixel) is. A
sentence is turned into a list of numbers before a language model reads it. A
house is described by numbers such as its floor area, its number of bedrooms
and its age. A machine learning \emph{model} is a recipe that turns input
numbers into output numbers (``this picture is a cat with probability
$0.97$''), and the recipe contains adjustable numbers, called
\emph{parameters}, that are tuned during training.

So before anything else we need to be fluent with numbers: what kinds of
numbers exist, how to calculate with them, and, crucially, how to reason
about numbers we do not know yet. That last skill is called \emph{algebra}. It
is the language in which every formula of machine learning is written.
\end{motivation}

\begin{objectives}
\begin{itemize}
  \item distinguish natural numbers, integers, rational numbers and real numbers,
    and use set and interval notation;
  \item calculate confidently with negative numbers, fractions, decimals and
    percentages;
  \item use exponents, roots and scientific notation (such as $3\times10^{-4}$);
  \item evaluate expressions using the correct order of operations;
  \item read and write summation notation $\sum$, the most common symbol in
    machine learning formulas;
  \item solve linear equations, simple systems of equations, and linear
    inequalities;
  \item work with absolute values and interpret them as distances.
\end{itemize}
\end{objectives}

%----------------------------------------------------------------------
\section{Kinds of numbers}\label{sec:kinds-of-numbers}
%----------------------------------------------------------------------

Numbers were invented to answer questions. ``How many sheep do I have?'' needs
one kind of number; ``how much do I owe?'' needs another; ``how long is the
diagonal of this square?'' needs yet another. Each new kind of question
enlarged the family of numbers. We will meet the family members in the order
in which they are usually discovered.

\subsection{Natural numbers}

The numbers we use for counting, $1, 2, 3, 4, \dots$, are called the
\term[natural number]{natural numbers}. The three dots mean ``and so on
forever'': there is no largest natural number, because you can always add
one more.

\begin{notation}
A collection of objects is called a \term{set}, and the objects in it are its
\term[element of a set]{elements}. We list the elements of a set between curly
braces. The set of natural numbers has a special symbol, a ``blackboard
bold'' letter N:
\[
  \N = \{1, 2, 3, 4, \dots\}.
\]
We write $5 \in \N$, read ``$5$ \emph{is an element of} $\N$'' or simply ``$5$
is in $\N$''. We write $-2 \notin \N$, read ``$-2$ is not in $\N$''.
\end{notation}

\begin{remark}
Some books include $0$ in the natural numbers and some do not. Nothing deep
depends on this choice. In these notes, $\N$ starts at $1$.
\end{remark}

\subsection{Integers}

Counting numbers cannot describe a temperature of five degrees below zero, a
debt of \$20, or a position $3$ metres below sea level. For that we need zero
and the negative numbers. The \term[integer]{integers} are the whole numbers,
positive, negative and zero:
\[
  \Z = \{\dots, -3, -2, -1, 0, 1, 2, 3, \dots\}.
\]
(The letter $\Z$ comes from the German word \emph{Zahlen}, ``numbers''.)

The best way to picture the integers is on a \term{number line}: a straight
line with $0$ in the middle, the positive numbers spaced evenly to the right
and the negative numbers to the left (\cref{fig:number-line}). Numbers further
to the right are larger. So $-1$ is larger than $-5$, even though $5$ is larger
than $1$: being $1$ degree below zero is warmer than being $5$ degrees below
zero.

\begin{figure}[htbp]
\centering
\begin{tikzpicture}[x=1.3cm]
  \draw[{Stealth[length=2.5mm]}-{Stealth[length=2.5mm]},thick] (-4.6,0) -- (4.6,0);
  \foreach \x in {-4,...,4}
    \draw (\x,0.08) -- (\x,-0.08) node[below=2pt,font=\small] {$\x$};
  \fill[PlotRed] (-2.5,0) circle (2.2pt) node[above=3pt,font=\small] {$-2.5$};
  \fill[PlotRed] (0.5,0) circle (2.2pt) node[above=3pt,font=\small] {$\tfrac12$};
  \fill[PlotRed] (1.41421,0) circle (2.2pt) node[above=3pt,font=\small] {$\sqrt2$};
  \fill[PlotRed] (3.14159,0) circle (2.2pt) node[above=3pt,font=\small] {$\pi$};
  \node[font=\small\itshape,text=MLremark] at (-2.5,-0.75) {$\longleftarrow$ smaller (negative)};
  \node[font=\small\itshape,text=MLremark] at (2.5,-0.75) {larger (positive) $\longrightarrow$};
\end{tikzpicture}
\caption{The number line. Every real number corresponds to exactly one point,
  and every point to exactly one real number. Integers are marked with ticks;
  a few non-integers are marked in red.}
\label{fig:number-line}
\end{figure}

\subsection{Rational numbers}

Sharing $3$ pizzas among $4$ people gives each person $\tfrac34$ of a pizza,
which is not a whole number. A \term{rational number} is any number that can
be written as a ratio (a fraction) of two integers,
\[
  \frac{p}{q}, \qquad\text{where $p$ and $q$ are integers and } q\neq 0.
\]
Examples are $\tfrac34$, $-\tfrac{5}{2}$, $\tfrac{7}{1}=7$ and
$0.25=\tfrac14$. The set of rational numbers is written $\Q$ (for
\emph{quotient}). Every integer is rational, because $n = \tfrac{n}{1}$.

When written as decimals, rational numbers either stop
($\tfrac14 = 0.25$) or repeat a pattern forever ($\tfrac13 = 0.333\ldots$,
$\tfrac{1}{7} = 0.142857142857\ldots$).

\subsection{Irrational and real numbers}

The ancient Greeks discovered, to their dismay, that some lengths are
\emph{not} ratios of integers. The diagonal of a square with sides of length
$1$ has length $\sqrt{2} = 1.41421356\ldots$, and it can be proved that no
fraction $p/q$ equals $\sqrt2$ exactly. Numbers like this, whose decimal
expansions never stop and never repeat, are called
\term[irrational number]{irrational}. Another famous example is
$\pi = 3.14159265\ldots$, the ratio of a circle's circumference to its
diameter.

Putting the rational and irrational numbers together gives the
\term[real number]{real numbers}, written $\R$. These are \emph{all} the points
on the number line, with no gaps. Almost all the numbers in machine learning
are real numbers: weights, probabilities, distances, losses and so on.

\begin{figure}[htbp]
\centering
\begin{tikzpicture}
  \draw[thick,MLdef,fill=MLdef!5] (0,0) ellipse (5.2cm and 2.4cm);
  \draw[thick,MLthm,fill=MLthm!6] (-0.8,0) ellipse (4.0cm and 1.85cm);
  \draw[thick,MLexample,fill=MLexample!7] (-1.6,0) ellipse (2.8cm and 1.3cm);
  \draw[thick,MLintuition,fill=MLintuition!12] (-2.3,0) ellipse (1.6cm and 0.8cm);
  \node[font=\large] at (-2.3,0.3) {$\N$};
  \node[font=\small] at (-2.3,-0.25) {$1,\ 2,\ 3,\ \dots$};
  \node[font=\large] at (0.25,0.45) {$\Z$};
  \node[font=\small] at (0.25,-0.2) {$0,\ -1,\ -7$};
  \node[font=\large] at (2.2,0.6) {$\Q$};
  \node[font=\small] at (2.2,-0.05) {$\tfrac12,\ -0.75$};
  \node[font=\large] at (4.25,0.6) {$\R$};
  \node[font=\small] at (4.25,-0.05) {$\sqrt2,\ \pi$};
\end{tikzpicture}
\caption{The number systems sit inside one another:
  $\N \subset \Z \subset \Q \subset \R$. Each ring shows a few numbers that
  belong to that set but not to the smaller ones inside it.}
\label{fig:nested-sets}
\end{figure}

\begin{notation}
If every element of a set $A$ is also an element of a set $B$, we say $A$ is a
\term{subset} of $B$ and write $A \subset B$. \Cref{fig:nested-sets} shows that
\[
  \N \subset \Z \subset \Q \subset \R .
\]
We can also describe a set by a property its elements satisfy, using
\term{set-builder notation}:
\[
  \{x \in \R : x > 0\}
\]
is read ``the set of all real numbers $x$ \emph{such that} $x$ is greater than
$0$''. The colon (sometimes written as a vertical bar $\mid$) means ``such
that''.
\end{notation}

\subsection{Intervals}

We often need to talk about all the real numbers between two values, for
example ``all probabilities'', which lie between $0$ and $1$. Such a set is
called an \term{interval}. A square bracket means the endpoint is included; a
round bracket means it is excluded.

\begin{center}
\renewcommand{\arraystretch}{1.25}
\begin{tabular}{@{}lll@{}}
\toprule
\textbf{Notation} & \textbf{Meaning} & \textbf{Name} \\
\midrule
$[a,b]$ & all $x$ with $a \le x \le b$ & closed interval \\
$(a,b)$ & all $x$ with $a < x < b$ & open interval \\
$[a,b)$ & all $x$ with $a \le x < b$ & half-open interval \\
$(a,\infty)$ & all $x$ with $x > a$ & unbounded interval \\
$(-\infty,b]$ & all $x$ with $x \le b$ & unbounded interval \\
\bottomrule
\end{tabular}
\end{center}

The symbol $\infty$ (``infinity'') is not a number; $(a,\infty)$ simply means
that the interval goes on forever to the right, so an infinity symbol always
gets a round bracket. On a number line we draw a filled dot for an included
endpoint and an open dot for an excluded one (\cref{fig:intervals}).

\begin{figure}[htbp]
\centering
\begin{tikzpicture}[x=0.75cm,y=0.9cm]
  \foreach \row/\lab in {0/{$[-2,3)$},-1/{$(1,\infty)$},-2/{$(-\infty,0]$}}{
    \draw[{Stealth[length=2mm]}-{Stealth[length=2mm]},thin,gray] (-4.6,\row) -- (5.6,\row);
    \foreach \x in {-4,...,5}
      \draw[gray] (\x,\row+0.07) -- (\x,\row-0.07);
    \node[anchor=east,font=\small] at (-4.8,\row) {\lab};
  }
  \foreach \x in {-4,...,5}
    \node[below,font=\scriptsize,gray] at (\x,-2.08) {$\x$};
  % [-2,3)
  \draw[line width=2.5pt,PlotBlue] (-2,0) -- (3,0);
  \fill[PlotBlue] (-2,0) circle (3.5pt);
  \draw[PlotBlue,thick,fill=white] (3,0) circle (3.5pt);
  % (1,infinity)
  \draw[line width=2.5pt,PlotBlue,-{Stealth[length=3mm]}] (1,-1) -- (5.6,-1);
  \draw[PlotBlue,thick,fill=white] (1,-1) circle (3.5pt);
  % (-infinity,0]
  \draw[line width=2.5pt,PlotBlue,{Stealth[length=3mm]}-] (-4.6,-2) -- (0,-2);
  \fill[PlotBlue] (0,-2) circle (3.5pt);
\end{tikzpicture}
\caption{Three intervals drawn on the number line. Filled dots are included
  endpoints, open dots are excluded endpoints.}
\label{fig:intervals}
\end{figure}

\begin{mlnote}[What numbers look like inside a computer]
A grayscale image is stored as a grid of integers from $0$ (black) to $255$
(white). Before feeding it to a neural network we usually divide every value
by $255$, so that all inputs lie in the interval $[0,1]$. Class labels such as
``cat'', ``dog'', ``bird'' are stored as integers $0,1,2$. Everything else
(weights, predictions, probabilities, losses) is a real number.

A computer cannot store infinitely many decimal digits, so it stores real
numbers \emph{approximately}, using a format called \emph{floating point}. The
common 32-bit format keeps about $7$ significant decimal digits. This is
usually plenty, but it means that extremely small numbers can be rounded to
$0$. We will see in \cref{ch:functions} that this is one reason why machine
learning works with \emph{logarithms} of probabilities.
\end{mlnote}

\sectionexercises

\begin{exercise}
For each number, list every set among $\N$, $\Z$, $\Q$, $\R$ that contains it:
$7$, $\ -3$, $\ 0$, $\ 0.25$, $\ -\tfrac23$, $\ \sqrt9$, $\ \sqrt5$, $\ \pi$.
\end{exercise}
\begin{solution}
$7$: $\N,\Z,\Q,\R$. \quad $-3$: $\Z,\Q,\R$. \quad $0$: $\Z,\Q,\R$ (not $\N$
with our convention). \quad $0.25=\tfrac14$: $\Q,\R$. \quad $-\tfrac23$:
$\Q,\R$. \quad $\sqrt9 = 3$: $\N,\Z,\Q,\R$. \quad $\sqrt5$: only $\R$ (it is
irrational). \quad $\pi$: only $\R$.
\end{solution}

\begin{exercise}
Write each set in interval notation and sketch it on a number line:
all $x$ with $-2 \le x < 5$; all $x$ with $x > 3$; all $x$ with $x \le 0$; all
$x$ with $0 \le x \le 1$.
\end{exercise}
\begin{solution}
$[-2,5)$ (filled dot at $-2$, open dot at $5$); $(3,\infty)$ (open dot at
$3$, arrow to the right); $(-\infty,0]$ (arrow to the left, filled dot at $0$);
$[0,1]$ (filled dots at both ends). The last one is the set of all possible
probabilities.
\end{solution}

\begin{exercise}
Show that $0.75$ and $0.272727\ldots$ (with ``$27$'' repeating forever) are
rational numbers by writing each as a fraction of integers. \emph{Hint for the
second: call the number $x$, compute $100x$, and subtract $x$.}
\end{exercise}
\begin{solution}
$0.75 = \tfrac{75}{100} = \tfrac34$. For the second, let
$x = 0.272727\ldots$. Then $100x = 27.272727\ldots$, and subtracting,
$100x - x = 27$, so $99x = 27$ and $x = \tfrac{27}{99} = \tfrac{3}{11}$.
(Check: $3 \div 11 = 0.2727\ldots$.)
\end{solution}

\begin{exercise}
A grayscale pixel is stored as an integer from $0$ to $255$.
How many different values can a pixel take? If we rescale by dividing by
$255$, in which interval do the rescaled values lie, and what is the rescaled
value of a pixel equal to $51$?
\end{exercise}
\begin{solution}
There are $256$ values ($0$ through $255$; do not forget $0$). The rescaled
values lie in $[0,1]$, since $0/255 = 0$ and $255/255 = 1$. A pixel value of
$51$ becomes $51/255 = 0.2$.
\end{solution}

%----------------------------------------------------------------------
\section{Arithmetic and its laws}\label{sec:arithmetic}
%----------------------------------------------------------------------

The four basic operations are addition ($+$), subtraction ($-$),
multiplication ($\times$) and division ($\div$ or $/$). The result of an
addition is called a \term{sum}, of a subtraction a \term{difference}, of a
multiplication a \term{product} and of a division a \term{quotient}. In
algebra the multiplication sign is often written as a dot, $3 \cdot 4$, or
left out entirely when letters are involved: $3x$ means $3 \times x$, and $ab$
means $a \times b$.

\subsection{The laws of arithmetic}

A handful of rules govern all calculations. They look obvious with small
numbers, but they are exactly what lets us rearrange complicated formulas
later on. In what follows, $a$, $b$ and $c$ stand for any real numbers.

\begin{theorem}[Laws of arithmetic]\label{thm:laws-arithmetic}
For all real numbers $a$, $b$, $c$:
\begin{enumerate}
  \item \textbf{Commutative laws} (order does not matter):
    $a + b = b + a$ and $ab = ba$.
  \item \textbf{Associative laws} (grouping does not matter):
    $(a+b)+c = a+(b+c)$ and $(ab)c = a(bc)$.
  \item \textbf{Distributive law} (multiplication spreads over addition):
    $a(b+c) = ab + ac$.
  \item \textbf{Identities}: $a + 0 = a$ and $a \cdot 1 = a$.
  \item \textbf{Inverses}: $a + (-a) = 0$, and if $a \neq 0$,
    $a \cdot \tfrac{1}{a} = 1$.
\end{enumerate}
\end{theorem}
\index{commutative law}\index{associative law}\index{distributive law}

The number $-a$ is the \term{additive inverse} (or negative) of $a$, and
$\tfrac1a$ is the \term{reciprocal} of $a$.

\begin{figure}[htbp]
\centering
\begin{tikzpicture}[scale=0.85]
  \draw[thick,fill=MLdef!15] (0,0) rectangle (4,3);
  \draw[thick,fill=MLexample!18] (4,0) rectangle (6,3);
  \node at (2,1.5) {$a\times b$};
  \node at (5,1.5) {$a\times c$};
  \draw[decorate,decoration={brace,amplitude=5pt,mirror}] (0,-0.1) -- (4,-0.1)
    node[midway,below=6pt] {$b$};
  \draw[decorate,decoration={brace,amplitude=5pt,mirror}] (4,-0.1) -- (6,-0.1)
    node[midway,below=6pt] {$c$};
  \draw[decorate,decoration={brace,amplitude=5pt}] (-0.1,0) -- (-0.1,3)
    node[midway,left=6pt] {$a$};
  \draw[decorate,decoration={brace,amplitude=5pt}] (0,3.1) -- (6,3.1)
    node[midway,above=6pt] {$b+c$};
\end{tikzpicture}
\caption{The distributive law as areas. The whole rectangle has area
  $a\times(b+c)$; its two pieces have areas $a\times b$ and $a\times c$.}
\label{fig:distributive}
\end{figure}

\begin{intuition}[The distributive law is about areas]
Think of $a \times (b+c)$ as the area of a rectangle of height $a$ and width
$b + c$ (\cref{fig:distributive}). Cutting it into two rectangles of widths $b$
and $c$ does not change the total area, which is therefore $ab + ac$. For
instance, $3 \times (4 + 2) = 3\times 6 = 18$ and $3\times4 + 3\times2 = 12 + 6
= 18$. The distributive law is also the trick behind mental arithmetic:
$7 \times 102 = 7\times100 + 7\times 2 = 714$.
\end{intuition}

\subsection{Negative numbers}

On the number line, adding a positive number moves you to the right and
subtracting a positive number moves you to the left. So $-3 + 5 = 2$ (start at
$-3$, move $5$ steps right) and $2 - 6 = -4$ (start at $2$, move $6$ steps
left).

Subtracting a negative number is the same as adding a positive one:
\[
  5 - (-3) = 5 + 3 = 8.
\]
One way to see this: subtraction asks ``how far is it from $-3$ up to $5$?'',
and on the number line that distance is $8$. Another way: taking away a debt of
\$3 leaves you \$3 richer.

For multiplication and division the rules of signs are:
\begin{center}
\renewcommand{\arraystretch}{1.2}
\begin{tabular}{@{}ll@{}}
\toprule
positive $\times$ positive $=$ positive & \quad $3 \times 4 = 12$\\
positive $\times$ negative $=$ negative & \quad $3 \times (-4) = -12$\\
negative $\times$ negative $=$ positive & \quad $(-3) \times (-4) = 12$\\
\bottomrule
\end{tabular}
\end{center}
and exactly the same rules hold for division, e.g.\ $(-12) \div (-4) = 3$.

\begin{intuition}[Why is negative times negative positive?]
Look at the pattern as the first number decreases by $1$ each time:
\[
  3 \times (-2) = -6,\quad 2 \times (-2) = -4,\quad 1\times(-2) = -2,\quad
  0 \times (-2) = 0,\quad (-1)\times(-2) = \ ?
\]
Each answer is $2$ larger than the one before, so the pattern forces
$(-1)\times(-2) = 2$. The rule ``negative times negative is positive'' is the
only rule that keeps the distributive law true; any other choice would break
arithmetic.
\end{intuition}

\subsection{Division by zero}

Division undoes multiplication: $12 \div 3 = 4$ because $3 \times 4 = 12$. So
$6 \div 0$ would have to be a number that gives $6$ when multiplied by $0$. But
anything times $0$ is $0$, so no such number exists.

\begin{pitfall}[Never divide by zero]
Division by zero is \emph{undefined}: expressions such as $\tfrac{5}{0}$ have
no meaning. (On the other hand, $\tfrac{0}{5} = 0$ is perfectly fine.) In
code, dividing by zero produces an error or a special ``not a number'' value,
and it is a common source of bugs in machine learning, for example when
normalising data whose spread happens to be zero. Practitioners therefore often
add a tiny number such as $10^{-8}$ to a denominator to keep it away from zero.
\end{pitfall}

\begin{pitfall}[Subtraction and division have no commutative or associative law]
$5 - 3 \neq 3 - 5$ and $8 \div 2 \neq 2 \div 8$. Grouping matters too:
$(10 - 5) - 2 = 3$, but $10 - (5 - 2) = 7$. If you want to rearrange
subtractions freely, rewrite them as additions of negatives:
$10 - 5 - 2 = 10 + (-5) + (-2)$, and now any order gives $3$.
\end{pitfall}

\begin{mlnote}[Linearity]
The distributive law $a(b + c) = ab + ac$ is the simplest example of a property
called \emph{linearity}: scaling a sum is the same as summing the scaled parts.
Linearity is the central idea of linear algebra (\cref{ch:linalg1}), and
linear algebra is the engine that runs every neural network.
\end{mlnote}

\sectionexercises

\begin{exercise}
Compute: (a) $-7 + 12$; (b) $4 - 9$; (c) $-3 - (-8)$; (d) $(-6) \times 4$;
(e) $(-5)\times(-7)$; (f) $(-36) \div (-4)$; (g) $0 \div 5$.
\end{exercise}
\begin{solution}
(a) $5$; (b) $-5$; (c) $-3 + 8 = 5$; (d) $-24$; (e) $35$; (f) $9$; (g) $0$.
\end{solution}

\begin{exercise}
Use the distributive law to compute $7 \times 98$ and $25 \times 44$ in your
head.
\end{exercise}
\begin{solution}
$7 \times 98 = 7 \times (100 - 2) = 700 - 14 = 686$.
$25 \times 44 = 25 \times (40 + 4) = 1000 + 100 = 1100$.
\end{solution}

\begin{exercise}
Give a numerical example showing that division is not associative, that is,
that $(a \div b) \div c$ can differ from $a \div (b \div c)$.
\end{exercise}
\begin{solution}
Take $a = 16$, $b = 4$, $c = 2$: $(16 \div 4) \div 2 = 4 \div 2 = 2$, while
$16 \div (4 \div 2) = 16 \div 2 = 8$.
\end{solution}

\begin{exercise}
A friend claims that $\tfrac{0}{0} = 1$ ``because anything divided by itself is
$1$''. Explain what is wrong.
\end{exercise}
\begin{solution}
``$\tfrac{x}{x} = 1$'' holds only for $x \neq 0$. If $\tfrac00$ were a number
$q$, we would need $0 \times q = 0$; but every number $q$ satisfies this, so
there is no single answer. For this reason $\tfrac00$ is left undefined (it is
called \emph{indeterminate}).
\end{solution}

%----------------------------------------------------------------------
\section{Fractions, decimals and percentages}\label{sec:fractions}
%----------------------------------------------------------------------

\subsection{What a fraction means}

A \term{fraction} $\tfrac{a}{b}$ has a \term{numerator} $a$ (on top) and a
\term{denominator} $b$ (on the bottom, never zero). The denominator says
\emph{into how many equal parts} the whole is cut; the numerator says
\emph{how many} of those parts we take. So $\tfrac34$ of a pizza means: cut the
pizza into $4$ equal slices and take $3$ of them. A fraction is also simply a
division: $\tfrac34 = 3 \div 4 = 0.75$.

\subsection{Equivalent fractions and simplifying}

Cutting each slice in half gives $8$ slices, of which we now have $6$:
$\tfrac34 = \tfrac68$. In general, multiplying or dividing the numerator
\emph{and} the denominator by the same non-zero number does not change the
value of a fraction:
\[
  \frac{a}{b} = \frac{a \times k}{b \times k} \qquad (k \neq 0).
\]
To \emph{simplify} a fraction, divide top and bottom by any common factor until
none is left: $\tfrac{18}{24} = \tfrac{18 \div 6}{24 \div 6} = \tfrac34$.

\subsection{Multiplying and dividing fractions}

To multiply fractions, multiply the numerators and multiply the
denominators:
\[
  \frac{a}{b}\times\frac{c}{d} = \frac{ac}{bd}.
\]
The word ``of'' usually signals multiplication: half of a third is
$\tfrac12 \times \tfrac13 = \tfrac16$.

To divide by a fraction, multiply by its reciprocal (flip it upside down):
\[
  \frac{a}{b} \div \frac{c}{d} = \frac{a}{b}\times\frac{d}{c} = \frac{ad}{bc}.
\]
Why? Ask how many quarters fit into $3$: each whole contains $4$ quarters, so
$3 \div \tfrac14 = 3 \times 4 = 12$.

\subsection{Adding and subtracting fractions}

You can only add fractions directly if they count slices of the same size,
that is, if they have the same denominator: $\tfrac{2}{7} + \tfrac{3}{7} =
\tfrac57$. Otherwise, first rewrite them with a \emph{common denominator}.

\begin{example}[Adding fractions]\label{ex:add-fractions}
Compute $\tfrac34 + \tfrac56$.

\sol A common denominator is $12$, since both $4$ and $6$ divide $12$. Then
$\tfrac34 = \tfrac{9}{12}$ and $\tfrac56 = \tfrac{10}{12}$, so
\[
  \frac34 + \frac56 = \frac{9}{12} + \frac{10}{12} = \frac{19}{12}.
\]
As a check, $0.75 + 0.8333\ldots = 1.5833\ldots$, and $19 \div 12 =
1.5833\ldots$.
\end{example}

\begin{pitfall}[Adding tops and bottoms]
It is tempting to write $\tfrac12 + \tfrac13 = \tfrac25$. This is wrong: the
answer must be bigger than $\tfrac12$, but $\tfrac25$ is smaller. The correct
calculation is $\tfrac12 + \tfrac13 = \tfrac36 + \tfrac26 = \tfrac56$.
\end{pitfall}

\begin{example}[Dividing fractions]
Compute $\tfrac23 \div \tfrac49$.

\sol Multiply by the reciprocal:
$\tfrac23 \times \tfrac94 = \tfrac{18}{12} = \tfrac32$.
\end{example}

\subsection{Decimals}

Decimals are fractions whose denominators are powers of ten:
$0.7 = \tfrac{7}{10}$, $0.25 = \tfrac{25}{100}$, $0.125 = \tfrac{125}{1000}$.
To turn a fraction into a decimal, divide: $\tfrac38 = 3 \div 8 = 0.375$.

\subsection{Percentages}

``Per cent'' means ``per hundred'', so a \term{percentage} is a fraction with
denominator $100$: $35\% = \tfrac{35}{100} = 0.35$. To find a percentage of a
quantity, multiply: $18\%$ of $250$ is $0.18 \times 250 = 45$.

The \term{percentage change} from an old value to a new value is
\[
  \text{percentage change} = \frac{\text{new} - \text{old}}{\text{old}}
  \times 100\%.
\]

\begin{example}[Percentage changes are not symmetric]\label{ex:percent-change}
A price rises from $80$ to $100$ and then falls back from $100$ to $80$. What
are the two percentage changes?

\sol Rise: $\tfrac{100 - 80}{80} \times 100\% = \tfrac{20}{80}\times100\% =
25\%$. Fall: $\tfrac{80 - 100}{100}\times100\% = -20\%$. The same absolute
change of $20$ is a $25\%$ rise but only a $20\%$ fall, because the
percentages are computed relative to different starting values.
\end{example}

\begin{pitfall}[Percent versus percentage points]
Suppose a classifier's accuracy improves from $90\%$ to $95\%$. The accuracy
has risen by $5$ \term[percentage point]{percentage points}, which is a
\emph{relative} increase of $\tfrac{5}{90} \approx 5.6\%$. Meanwhile the
\emph{error rate} has fallen from $10\%$ to $5\%$, a relative reduction of
$50\%$: the model makes half as many mistakes. All three statements describe
the same improvement. When you read ``a 50\% improvement'' in a paper, always
ask: $50\%$ of \emph{what}?
\end{pitfall}

\subsection{Ratios}

A \term{ratio} compares two quantities. If a data set contains $750$ normal
emails and $250$ spam emails, the ratio of normal to spam is $750 : 250$, which
simplifies to $3 : 1$. Out of every $4$ emails, $3$ are normal and $1$ is spam,
so the \emph{fraction} of spam is $\tfrac14 = 25\%$.

\begin{mlnote}
Fractions and percentages are everywhere in machine learning: accuracy
($\tfrac{\text{correct predictions}}{\text{all predictions}}$), train/test
splits (``use $80\%$ of the data for training''), class imbalance (``spam is
$1$ in $4$ emails'') and dropout rates (``switch off $10\%$ of the neurons'').
Being precise about \emph{what} a percentage is a percentage \emph{of} avoids
many misunderstandings.
\end{mlnote}

\sectionexercises

\begin{exercise}
Simplify: (a) $\tfrac{18}{24}$; (b) $\tfrac{35}{49}$; (c) $\tfrac{120}{36}$.
\end{exercise}
\begin{solution}
(a) $\tfrac34$ (divide by $6$); (b) $\tfrac57$ (divide by $7$);
(c) $\tfrac{10}{3}$ (divide by $12$).
\end{solution}

\begin{exercise}
Compute and simplify: (a) $\tfrac25 + \tfrac13$; (b) $\tfrac78 - \tfrac16$;
(c) $\tfrac34 \times \tfrac89$; (d) $\tfrac56 \div \tfrac{10}{3}$.
\end{exercise}
\begin{solution}
(a) $\tfrac{6}{15} + \tfrac{5}{15} = \tfrac{11}{15}$.
(b) $\tfrac{21}{24} - \tfrac{4}{24} = \tfrac{17}{24}$.
(c) $\tfrac{24}{36} = \tfrac23$.
(d) $\tfrac56 \times \tfrac{3}{10} = \tfrac{15}{60} = \tfrac14$.
\end{solution}

\begin{exercise}
(a) What is $15\%$ of $80$? (b) $45$ is what percentage of $180$? (c) A
price rises from $40$ to $50$; what is the percentage change? (d) It then
falls back to $40$; what is the percentage change now?
\end{exercise}
\begin{solution}
(a) $0.15 \times 80 = 12$. (b) $\tfrac{45}{180} = 0.25 = 25\%$.
(c) $\tfrac{50-40}{40} = 25\%$. (d) $\tfrac{40-50}{50} = -20\%$.
\end{solution}

\begin{exercise}
Model A has an accuracy of $80\%$ and model B an accuracy of $84\%$.
(a) By how many percentage points is B better? (b) What is the relative
increase in accuracy? (c) What is the relative reduction in the error rate?
\end{exercise}
\begin{solution}
(a) $4$ percentage points. (b) $\tfrac{84 - 80}{80} = 5\%$. (c) The error rate
falls from $20\%$ to $16\%$, a relative reduction of $\tfrac{20-16}{20} =
20\%$.
\end{solution}

\begin{exercise}
A data set of $1000$ emails has a $3:1$ ratio of normal emails to spam. How
many spam emails are there? If we keep $80\%$ of the data for training and
the split preserves the ratio, how many spam emails are in the training set?
\end{exercise}
\begin{solution}
Spam is $1$ part out of $3+1=4$, so there are $\tfrac14 \times 1000 = 250$ spam
emails. The training set has $0.8 \times 1000 = 800$ emails, of which
$\tfrac14 \times 800 = 200$ are spam.
\end{solution}

%----------------------------------------------------------------------
\section{Exponents, roots and scientific notation}\label{sec:exponents}
%----------------------------------------------------------------------

\subsection{Powers}

Repeated addition is multiplication: $3 + 3 + 3 + 3 = 4 \times 3$. Repeated
multiplication is \emph{exponentiation}:
\[
  3^4 = 3 \times 3 \times 3 \times 3 = 81.
\]
In $a^n$, the number $a$ is the \term{base} and $n$ is the
\term{exponent} (or \term{power}); we read $a^n$ as ``$a$ to the power $n$''.
Two powers have their own names: $a^2$ is ``$a$ \emph{squared}'' (the area of
a square with side $a$) and $a^3$ is ``$a$ \emph{cubed}'' (the volume of a cube
with side $a$).

\subsection{The laws of exponents}

Counting factors explains every rule. For example,
$2^3 \times 2^2 = (2\times2\times2)\times(2\times2) = 2^5$: three factors plus
two factors make five factors.

\begin{theorem}[Laws of exponents]\label{thm:exponent-laws}
For real numbers $a, b \neq 0$ and exponents $m, n$:
\begin{align*}
  a^m \cdot a^n &= a^{m+n}, &
  \frac{a^m}{a^n} &= a^{m-n}, &
  (a^m)^n &= a^{mn}, \\
  (ab)^n &= a^n b^n, &
  \left(\frac{a}{b}\right)^n &= \frac{a^n}{b^n}, &
  a^0 &= 1, \qquad a^{-n} = \frac{1}{a^n}.
\end{align*}
\end{theorem}

The last two rules may look strange: what does it mean to multiply $2$ by
itself zero times, or minus three times? They are \emph{definitions}, chosen so
that the other rules keep working. Watch the pattern as we repeatedly divide
by $2$:
\[
  2^3 = 8,\quad 2^2 = 4,\quad 2^1 = 2,\quad 2^0 = 1,\quad
  2^{-1} = \tfrac12,\quad 2^{-2} = \tfrac14,\quad 2^{-3} = \tfrac18 .
\]
Each step down in the exponent divides by $2$. Continuing the pattern forces
$2^0 = 1$ and $2^{-n} = \tfrac{1}{2^n}$.

\subsection{Roots and fractional exponents}

The \term{square root} of a number $a \ge 0$, written $\sqrt a$, is the
non-negative number whose square is $a$: $\sqrt{9} = 3$ because $3^2 = 9$.
(Note that $(-3)^2 = 9$ as well, but the symbol $\sqrt{\ }$ always means the
non-negative root.) Similarly the \emph{cube root} $\sqrt[3]{a}$ is the number
whose cube is $a$: $\sqrt[3]{8} = 2$.

Roots are powers with fractional exponents. Since
$(a^{1/2})^2 = a^{(1/2)\cdot 2} = a^1 = a$ by the power rule, the number
$a^{1/2}$ must be $\sqrt a$. In general,
\[
  a^{1/n} = \sqrt[n]{a}, \qquad a^{m/n} = \left(\sqrt[n]{a}\right)^m .
\]
For example, $8^{2/3} = (\sqrt[3]{8})^2 = 2^2 = 4$, and
$27^{-1/3} = \tfrac{1}{\sqrt[3]{27}} = \tfrac13$.

\begin{pitfall}[Powers do not distribute over sums]
$(a + b)^2$ is \textbf{not} $a^2 + b^2$. With numbers:
$(3 + 4)^2 = 7^2 = 49$, but $3^2 + 4^2 = 9 + 16 = 25$. The correct expansion,
$(a+b)^2 = a^2 + 2ab + b^2$, is pictured in \cref{fig:square-of-sum}: the
missing $24 = 2\times 3\times 4$ is the area of the two rectangles. Likewise
$\sqrt{a + b} \neq \sqrt a + \sqrt b$: $\sqrt{9 + 16} = 5$, but
$\sqrt9 + \sqrt{16} = 7$. Powers and roots \emph{do} distribute over products:
$(ab)^2 = a^2b^2$ and $\sqrt{ab} = \sqrt a\sqrt b$ for $a, b \ge 0$.
\end{pitfall}

\begin{figure}[htbp]
\centering
\begin{tikzpicture}[scale=0.75]
  \draw[thick,fill=MLdef!15] (0,0) rectangle (3,3);
  \draw[thick,fill=MLexample!18] (3,0) rectangle (4.8,3);
  \draw[thick,fill=MLexample!18] (0,3) rectangle (3,4.8);
  \draw[thick,fill=MLintuition!22] (3,3) rectangle (4.8,4.8);
  \node at (1.5,1.5) {$a^2$};
  \node at (3.9,1.5) {$ab$};
  \node at (1.5,3.9) {$ab$};
  \node at (3.9,3.9) {$b^2$};
  \draw[decorate,decoration={brace,amplitude=5pt,mirror}] (0,-0.1) -- (3,-0.1)
    node[midway,below=6pt] {$a$};
  \draw[decorate,decoration={brace,amplitude=5pt,mirror}] (3,-0.1) -- (4.8,-0.1)
    node[midway,below=6pt] {$b$};
  \draw[decorate,decoration={brace,amplitude=5pt}] (-0.1,0) -- (-0.1,3)
    node[midway,left=6pt] {$a$};
  \draw[decorate,decoration={brace,amplitude=5pt}] (-0.1,3) -- (-0.1,4.8)
    node[midway,left=6pt] {$b$};
\end{tikzpicture}
\caption{A square with side $a+b$ has area $(a+b)^2 = a^2 + 2ab + b^2$:
  besides the squares $a^2$ and $b^2$ there are two rectangles of area $ab$.}
\label{fig:square-of-sum}
\end{figure}

\begin{pitfall}[Negative signs and powers]
$-3^2$ means $-(3^2) = -9$, because the power applies only to the $3$. If you
want to square $-3$, write $(-3)^2 = 9$. Also, the square root of a negative
number, such as $\sqrt{-4}$, is not a real number (no real number squared is
negative). Finally, $\sqrt{x^2} = \abs{x}$, not $x$: for example
$\sqrt{(-5)^2} = \sqrt{25} = 5$. (The absolute value $\abs{x}$ is explained in
\cref{sec:absolute-value}.)
\end{pitfall}

\subsection{Scientific notation}

Machine learning involves both very large and very small numbers. Writing them
out in full is error-prone, so we use \term{scientific notation}: a number
between $1$ and $10$ multiplied by a power of $10$.
\[
  0.0003 = 3 \times 10^{-4}, \qquad
  175\,000\,000\,000 = 1.75 \times 10^{11}.
\]
The exponent counts how many places the decimal point moves: right for
positive exponents, left for negative ones. To multiply numbers in scientific
notation, multiply the leading parts and add the exponents:
\[
  (3 \times 10^{-4}) \times (2 \times 10^{6}) = 6 \times 10^{2} = 600 .
\]
In computer code, $3 \times 10^{-4}$ is usually written \texttt{3e-4}.

\begin{mlnote}
Typical numbers in deep learning: a learning rate of \texttt{3e-4}, a weight
decay of \texttt{1e-2}, a small ``safety'' constant $\epsilon$ of \texttt{1e-8}
added to denominators, and models with $1.75\times10^{11}$ parameters (the
size of the GPT-3 language model). Being able to read and compare such numbers
at a glance is a basic practical skill.
\end{mlnote}

\sectionexercises

\begin{exercise}
Evaluate: (a) $2^3\times2^4$; (b) $\tfrac{5^6}{5^4}$; (c) $(3^2)^3$;
(d) $10^0$; (e) $4^{-2}$; (f) $9^{1/2}$; (g) $8^{2/3}$; (h) $16^{-3/4}$.
\end{exercise}
\begin{solution}
(a) $2^7 = 128$; (b) $5^2 = 25$; (c) $3^6 = 729$; (d) $1$; (e)
$\tfrac{1}{16}$; (f) $3$; (g) $(\sqrt[3]8)^2 = 4$; (h)
$\tfrac{1}{(\sqrt[4]{16})^3} = \tfrac{1}{2^3} = \tfrac18$.
\end{solution}

\begin{exercise}
Simplify, assuming $x \neq 0$: (a) $\tfrac{x^3 \cdot x^5}{x^2}$; (b) $(2x^2)^3$;
(c) $\sqrt{x^4}$; (d) $\tfrac{(x^{-1})^2}{x^{-3}}$.
\end{exercise}
\begin{solution}
(a) $x^{3+5-2} = x^6$. (b) $2^3 (x^2)^3 = 8x^6$. (c) $x^2$ (which is never
negative, so no absolute value is needed). (d) $x^{-2}/x^{-3} = x^{-2-(-3)} =
x^1 = x$.
\end{solution}

\begin{exercise}
True or false? Give a counterexample for each false statement.
(a) $(a + b)^2 = a^2 + b^2$; (b) $\sqrt{ab} = \sqrt a\sqrt b$ for $a,b\ge0$;
(c) $\sqrt{a+b} = \sqrt a + \sqrt b$; (d) $(ab)^3 = a^3b^3$;
(e) $2^{a+b} = 2^a + 2^b$.
\end{exercise}
\begin{solution}
(a) False: $a=b=1$ gives $4$ versus $2$. (b) True. (c) False: $a=9$, $b=16$
gives $5$ versus $7$. (d) True. (e) False: $a = b = 1$ gives $2^2 = 4$ versus
$2 + 2 = 4$, which happens to agree, but $a=1$, $b=2$ gives $2^3 = 8$ versus
$2 + 4 = 6$. (The correct rule is $2^{a+b} = 2^a \cdot 2^b$.)
\end{solution}

\begin{exercise}
Write $0.00005$ and $3\,400\,000$ in scientific notation, and compute
$(4 \times 10^{3}) \times (2.5\times 10^{-5})$.
\end{exercise}
\begin{solution}
$0.00005 = 5\times10^{-5}$; $3\,400\,000 = 3.4\times10^{6}$;
$(4\times2.5)\times10^{3-5} = 10\times10^{-2} = 10^{-1} = 0.1$.
\end{solution}

\begin{exercise}
A language model has $7 \times 10^9$ parameters, each stored using $2$ bytes.
How many bytes are needed to store the model? Express the answer in gigabytes,
where $1$ gigabyte $= 10^9$ bytes.
\end{exercise}
\begin{solution}
$7\times10^9 \times 2 = 1.4\times10^{10}$ bytes, which is
$\tfrac{1.4\times10^{10}}{10^9} = 14$ gigabytes.
\end{solution}

%----------------------------------------------------------------------
\section{Order of operations}\label{sec:order-of-operations}
%----------------------------------------------------------------------

What is $2 + 3 \times 4$? If you add first you get $20$; if you multiply first
you get $14$. Mathematics needs a single agreed answer, and the convention is
that multiplication comes first, so $2 + 3\times 4 = 14$. The full convention,
the \term{order of operations}, is:

\begin{enumerate}[label=\arabic*.]
  \item \textbf{Parentheses} (brackets): work out anything inside brackets
    first, innermost brackets first.
  \item \textbf{Exponents} (powers and roots).
  \item \textbf{Multiplication and division}, from left to right.
  \item \textbf{Addition and subtraction}, from left to right.
\end{enumerate}

Many people remember this as PEMDAS or BODMAS. Note that multiplication and
division have \emph{equal} rank and are done left to right, and likewise for
addition and subtraction.

\begin{example}[Order of operations]
Evaluate (a) $8 - 2 \times 3 + 4$; (b) $20 \div 4 \times 5$;
(c) $2 \times 3^2$; (d) $(2+3)^2 - 4 \times 5$;
(e) $3 + 12 \div (5 - 2) \times 2$.

\sol
(a) Multiply first: $8 - 6 + 4$; then left to right: $2 + 4 = 6$.
(b) Left to right: $20 \div 4 = 5$, then $5 \times 5 = 25$ (not
$20 \div 20 = 1$).
(c) Power first: $2 \times 9 = 18$ (not $6^2 = 36$).
(d) Brackets, then power: $5^2 - 4\times5 = 25 - 20 = 5$.
(e) Brackets: $3 + 12 \div 3 \times 2$; then left to right:
$12\div3 = 4$, $4\times2 = 8$; finally $3 + 8 = 11$.
\end{example}

A fraction bar groups like brackets: $\dfrac{6 + 4}{2}$ means $(6+4) \div 2 =
5$, not $6 + 4 \div 2 = 8$.

\begin{pitfall}[Ambiguous expressions]
Expressions such as $1/2x$ are ambiguous: is it $\tfrac{1}{2x}$ or
$\tfrac12 x$? Different people (and different calculators) disagree. When
writing, always add brackets to remove any doubt: $1/(2x)$ or $(1/2)x$.
\end{pitfall}

\begin{remark}[Order of operations in code]
Programming languages follow the same conventions. In Python, powers are
written \texttt{**}, and \texttt{-2**2} evaluates to \texttt{-4}, exactly like
$-2^2$ in mathematics, because the power is applied before the minus sign.
\end{remark}

\sectionexercises

\begin{exercise}
Evaluate: (a) $6 + 4\times3 - 2$; (b) $(6+4)\times(3-2)$; (c) $48\div6\div2$;
(d) $2\times5^2 - 3^2$; (e) $-3^2 + (-3)^2$; (f) $(8 - 2\times3)^3$.
\end{exercise}
\begin{solution}
(a) $6 + 12 - 2 = 16$. (b) $10 \times 1 = 10$. (c) $8 \div 2 = 4$.
(d) $50 - 9 = 41$. (e) $-9 + 9 = 0$. (f) $(8-6)^3 = 2^3 = 8$.
\end{solution}

\begin{exercise}
Insert brackets into $2 + 3 \times 4 - 1$ so that its value is (a) $19$;
(b) $11$; (c) $15$.
\end{exercise}
\begin{solution}
(a) $(2+3)\times4 - 1 = 19$. (b) $2 + 3\times(4-1) = 11$.
(c) $(2+3)\times(4-1) = 15$.
\end{solution}

\begin{exercise}
Evaluate $\dfrac{2^3 + 4}{6 - 2\times 3 + 3}$ and $\sqrt{3^2 + 4^2}$.
\end{exercise}
\begin{solution}
Numerator $8 + 4 = 12$, denominator $6 - 6 + 3 = 3$, so the value is $4$.
And $\sqrt{9 + 16} = \sqrt{25} = 5$.
\end{solution}

%----------------------------------------------------------------------
\section{Variables, expressions and algebraic thinking}\label{sec:algebra}
%----------------------------------------------------------------------

\subsection{Variables}

Algebra begins with a simple but powerful idea: use a letter to stand for a
number. The letter is called a \term{variable}. It may stand for a number we
do not know yet (``find $x$''), or for a number that can change (``the price
depends on the size $x$ of the house''), or for \emph{any} number at all (``for
every $x$, $x + 0 = x$'').

An \term{expression} is a combination of numbers, variables and operations,
such as $3x + 2$ or $x^2 - 5x + 6$. \emph{Evaluating} an expression means
substituting a value for each variable and calculating. For $x = 4$,
$3x + 2 = 3 \times 4 + 2 = 14$.

\begin{intuition}[A variable is a labelled box]
Think of a variable as a box with a name written on it. The expression
$3x + 2$ is a recipe: ``take whatever is in box $x$, multiply it by $3$, then
add $2$''. The recipe makes sense before we know what is in the box. This is
exactly how a machine learning model is written down: as a recipe in terms of
inputs (boxes filled by the data) and parameters (boxes filled during
training).
\end{intuition}

\subsection{Terms, coefficients and like terms}

In $5x^2 - 3x + 7$, the pieces separated by $+$ or $-$ signs are called
\term[term of an expression]{terms}: $5x^2$, $-3x$ and $7$. The number in front
of a variable is its \term{coefficient} ($5$ and $-3$ here), and a term
without a variable is a \term{constant}. \term[like terms]{Like terms} contain the same
variables raised to the same powers, and they can be combined by adding their
coefficients:
\[
  3x + 5 + 2x - 1 = (3+2)x + (5 - 1) = 5x + 4 .
\]
This is just the distributive law in reverse: $3x + 2x = (3 + 2)x$. Unlike
terms, such as $x$ and $x^2$, cannot be combined.

\subsection{Expanding brackets}

Expanding means multiplying out brackets using the distributive law:
$3(x + 2) = 3x + 6$ and $-2(x - 5) = -2x + 10$ (careful with the signs). For a
product of two brackets, multiply every term in the first bracket by every
term in the second:
\[
  (x + 2)(x + 3) = x\cdot x + x \cdot 3 + 2 \cdot x + 2 \cdot 3
  = x^2 + 5x + 6.
\]
\Cref{fig:area-model} shows this as the area of a rectangle with sides $x + 2$
and $x + 3$.

\begin{figure}[htbp]
\centering
\begin{tikzpicture}[scale=0.7]
  \draw[thick,fill=MLdef!15] (0,0) rectangle (4,4);
  \draw[thick,fill=MLexample!18] (4,0) rectangle (7,4);
  \draw[thick,fill=MLexample!18] (0,4) rectangle (4,6);
  \draw[thick,fill=MLintuition!22] (4,4) rectangle (7,6);
  \node at (2,2) {$x^2$};
  \node at (5.5,2) {$3x$};
  \node at (2,5) {$2x$};
  \node at (5.5,5) {$6$};
  \draw[decorate,decoration={brace,amplitude=5pt,mirror}] (0,-0.1) -- (4,-0.1)
    node[midway,below=6pt] {$x$};
  \draw[decorate,decoration={brace,amplitude=5pt,mirror}] (4,-0.1) -- (7,-0.1)
    node[midway,below=6pt] {$3$};
  \draw[decorate,decoration={brace,amplitude=5pt}] (-0.1,0) -- (-0.1,4)
    node[midway,left=6pt] {$x$};
  \draw[decorate,decoration={brace,amplitude=5pt}] (-0.1,4) -- (-0.1,6)
    node[midway,left=6pt] {$2$};
\end{tikzpicture}
\caption{The area model for $(x+2)(x+3)$. The four pieces have areas
  $x^2$, $3x$, $2x$ and $6$, giving $x^2 + 5x + 6$.}
\label{fig:area-model}
\end{figure}

Three products appear so often that they are worth memorising:
\begin{align*}
  (a+b)^2 &= a^2 + 2ab + b^2, \\
  (a-b)^2 &= a^2 - 2ab + b^2, \\
  (a+b)(a-b) &= a^2 - b^2 .
\end{align*}
The second identity will reappear many times, because machine learning is full
of squared differences such as $(y - \hat y)^2$.

\subsection{Factoring}

\term[factoring]{Factoring} is expanding in reverse: writing an expression as a product.
Pull out a common factor, $6x + 9 = 3(2x + 3)$, or find two brackets that
multiply to the expression, $x^2 + 5x + 6 = (x+2)(x+3)$. To factor
$x^2 + bx + c$, look for two numbers that \emph{add} to $b$ and
\emph{multiply} to $c$; here $2 + 3 = 5$ and $2 \times 3 = 6$. Factoring will
help us solve equations involving $x^2$ in \cref{ch:functions}.

\subsection{Formulas}

A \term{formula} is an equation that tells you how to compute one quantity from
others. Examples are the area of a rectangle, $A = \ell w$; the conversion from
degrees Celsius to degrees Fahrenheit, $F = \tfrac95 C + 32$; and the
prediction of a simple model of house prices,
\[
  \text{predicted price} = w \times \text{size} + b ,
\]
where $w$ and $b$ are numbers that the model learns from data.

\begin{mlnote}[Models are expressions with parameters]
In the house-price formula, ``size'' is an \emph{input} (it changes from house
to house) while $w$ and $b$ are \emph{parameters} (they are fixed once the
model has been trained). In machine learning, $w$ is called a \term{weight} and
$b$ a \term{bias}. A large neural network is the same idea on a huge scale: a
very long expression with millions or billions of weights. Training means
finding values of the weights that make the expression's outputs match the
data.
\end{mlnote}

\sectionexercises

\begin{exercise}
Evaluate $2x^2 - 3x + 1$ at $x = 3$, at $x = -1$ and at $x = \tfrac12$.
\end{exercise}
\begin{solution}
At $x=3$: $18 - 9 + 1 = 10$. At $x=-1$: $2 + 3 + 1 = 6$. At $x=\tfrac12$:
$\tfrac12 - \tfrac32 + 1 = 0$.
\end{solution}

\begin{exercise}
Simplify: (a) $4a + 3b - a + 5b$; (b) $2(x-3) - 3(x+1)$;
(c) $x(x + 2) - (x^2 - 4x)$.
\end{exercise}
\begin{solution}
(a) $3a + 8b$. (b) $2x - 6 - 3x - 3 = -x - 9$.
(c) $x^2 + 2x - x^2 + 4x = 6x$.
\end{solution}

\begin{exercise}
Expand: (a) $(x+4)(x-2)$; (b) $(2x - 3)^2$; (c) $(x+5)(x-5)$;
(d) $(a - b)(a^2 + ab + b^2)$.
\end{exercise}
\begin{solution}
(a) $x^2 - 2x + 4x - 8 = x^2 + 2x - 8$. (b) $4x^2 - 12x + 9$.
(c) $x^2 - 25$. (d) $a^3 + a^2b + ab^2 - a^2b - ab^2 - b^3 = a^3 - b^3$.
\end{solution}

\begin{exercise}
Factor: (a) $12x + 18$; (b) $x^2 + 7x + 12$; (c) $x^2 - 9$;
(d) $x^2 - x - 6$.
\end{exercise}
\begin{solution}
(a) $6(2x + 3)$. (b) $(x+3)(x+4)$, since $3+4=7$ and $3\times4=12$.
(c) $(x-3)(x+3)$. (d) $(x-3)(x+2)$, since $-3+2=-1$ and $(-3)\times2=-6$.
\end{solution}

\begin{exercise}
A model predicts house prices with the formula
$\text{price} = 3000 \times \text{size} + 50\,000$, with size in square
metres. What does it predict for a house of $80$ square metres? If the true
price is $300\,000$, what is the error $\text{true} - \text{predicted}$?
\end{exercise}
\begin{solution}
$3000 \times 80 + 50\,000 = 240\,000 + 50\,000 = 290\,000$. The error is
$300\,000 - 290\,000 = 10\,000$: the model under-predicts by $10\,000$.
\end{solution}

%----------------------------------------------------------------------
\section{Summation and product notation}\label{sec:summation}
%----------------------------------------------------------------------

\begin{motivation}
Machine learning learns from data sets with thousands or millions of
examples, and its formulas constantly say ``add this up over all the
examples''. Writing out $x_1 + x_2 + x_3 + \cdots + x_{1000000}$ is
impractical, so mathematicians use a compact shorthand: the capital Greek
letter sigma, $\sum$. You will see it in almost every machine learning formula,
so it is worth getting completely comfortable with it now.
\end{motivation}

\subsection{Subscripts}

When we have many related numbers, we give them the same letter and tell them
apart with a small number written below, called a \term{subscript} or
\term{index}. A data set of four measurements might be
\[
  x_1 = 2,\quad x_2 = 5,\quad x_3 = 3,\quad x_4 = 7 .
\]
We read $x_3$ as ``$x$ sub three'' or ``$x$ three''. A general element is
written $x_i$, where $i$ can be any of $1, 2, 3, 4$.

\subsection{Sigma notation}

\begin{definition}[Summation notation]\label{def:summation}
For numbers $x_1, x_2, \dots, x_n$,
\[
  \sum_{i=1}^{n} x_i = x_1 + x_2 + \cdots + x_n .
\]
This is read ``the sum of $x_i$ for $i$ from $1$ to $n$''. The variable $i$ is
the \term[index of summation]{index of summation}; it starts at the value
below the $\sum$ and increases in steps of $1$ up to the value on top.
\end{definition}
\index{summation notation}\index{sigma notation|see{summation notation}}

\begin{example}[Reading sums]\label{ex:sums}
Compute (a) $\sum_{i=1}^{4} i$; (b) $\sum_{i=1}^{3} i^2$;
(c) $\sum_{k=0}^{3} 2^k$; (d) $\sum_{i=1}^4 x_i$ for the data
$x_1 = 2, x_2 = 5, x_3 = 3, x_4 = 7$.

\sol
(a) $1 + 2 + 3 + 4 = 10$.
(b) $1^2 + 2^2 + 3^2 = 1 + 4 + 9 = 14$.
(c) $2^0 + 2^1 + 2^2 + 2^3 = 1 + 2 + 4 + 8 = 15$ (the index may start at $0$).
(d) $2 + 5 + 3 + 7 = 17$.
\end{example}

The name of the index does not matter: $\sum_{i=1}^4 x_i$, $\sum_{j=1}^4 x_j$
and $\sum_{k=1}^4 x_k$ all mean the same thing. For this reason the index is
called a \emph{dummy variable}.

\subsection{Rules for sums}

Because a sum is just repeated addition, the laws of arithmetic give us some
useful rules.

\begin{proposition}[Rules for sums]\label{prop:sum-rules}
For numbers $x_1,\dots,x_n$, $y_1, \dots, y_n$ and a constant $c$:
\begin{align*}
  \sum_{i=1}^{n} c\,x_i &= c\sum_{i=1}^{n} x_i
    && \text{(constants can be pulled out)},\\
  \sum_{i=1}^{n} (x_i + y_i) &= \sum_{i=1}^{n} x_i + \sum_{i=1}^{n} y_i
    && \text{(sums can be split)},\\
  \sum_{i=1}^{n} c &= n\,c
    && \text{(adding $c$ to itself $n$ times)}.
\end{align*}
\end{proposition}

\begin{proof}
The first rule is the distributive law:
$c x_1 + c x_2 + \dots + c x_n = c(x_1 + \dots + x_n)$. The second rule just
reorders the terms, which the commutative and associative laws allow. The
third rule is the definition of multiplication.
\end{proof}

\begin{pitfall}[What you cannot do with sums]
In general,
\[
  \sum_i x_i y_i \neq \Bigl(\sum_i x_i\Bigr)\Bigl(\sum_i y_i\Bigr)
  \qquad\text{and}\qquad
  \sum_i x_i^2 \neq \Bigl(\sum_i x_i\Bigr)^2 .
\]
For the data $2, 5, 3, 7$: $\sum_i x_i^2 = 4 + 25 + 9 + 49 = 87$, whereas
$\bigl(\sum_i x_i\bigr)^2 = 17^2 = 289$.
\end{pitfall}

\subsection{The mean}

The \term[mean]{mean} (average) of $n$ numbers is their sum divided by $n$. It
is written with a bar:
\[
  \bar x = \frac{1}{n}\sum_{i=1}^{n} x_i .
\]
For the data $2, 5, 3, 7$, $\bar x = \tfrac{17}{4} = 4.25$.

\subsection{Product notation}

The capital Greek letter pi, $\prod$, does for multiplication what $\sum$ does
for addition:
\[
  \prod_{i=1}^{n} x_i = x_1 \times x_2 \times \cdots \times x_n .
\]
For example $\prod_{i=1}^{4} i = 1\times2\times3\times4 = 24$. This
particular product is so common that it has its own name and symbol:
$n! = 1 \times 2 \times \cdots \times n$ is called ``$n$
\term{factorial}'', so $4! = 24$. By convention, $0! = 1$.

\begin{mlnote}[The mean squared error]
Suppose a model makes predictions $\hat y_1, \dots, \hat y_n$ (read ``$y$
hat'') for $n$ data points whose true values are $y_1, \dots, y_n$. The most
common way to measure how wrong the model is, is the
\term{mean squared error}:
\[
  \MSE = \frac{1}{n}\sum_{i=1}^{n} (y_i - \hat y_i)^2 .
\]
In words: for each data point, compute the error $y_i - \hat y_i$, square it
(so that positive and negative errors both count, and big errors count a lot),
and average over all data points. For true values $3, 5, 4, 6$ and predictions
$2, 5, 5, 8$, the errors are $1, 0, -1, -2$, their squares are $1, 0, 1, 4$,
and $\MSE = \tfrac{6}{4} = 1.5$. Training a model usually means making a
quantity like this as small as possible.

Products appear when we combine probabilities: the probability that a model
assigns to an entire data set is often a product $\prod_{i=1}^n p_i$ of the
probabilities of the individual examples (\cref{ch:statistics}).
\end{mlnote}

\sectionexercises

\begin{exercise}
Compute: (a) $\sum_{i=1}^{5}(2i - 1)$; (b) $\sum_{k=1}^{4} k^2$;
(c) $\sum_{j=0}^{3} 3$; (d) $\prod_{i=1}^{5} i$.
\end{exercise}
\begin{solution}
(a) $1 + 3 + 5 + 7 + 9 = 25$. (b) $1 + 4 + 9 + 16 = 30$.
(c) There are four terms ($j = 0,1,2,3$), each equal to $3$, so the sum is
$12$. (d) $5! = 120$.
\end{solution}

\begin{exercise}
Write in sigma notation: (a) $1 + 2 + \cdots + 100$;
(b) $x_1^2 + x_2^2 + \cdots + x_n^2$; (c) $w_1 x_1 + w_2 x_2 + w_3 x_3$.
\end{exercise}
\begin{solution}
(a) $\sum_{i=1}^{100} i$. (b) $\sum_{i=1}^{n} x_i^2$.
(c) $\sum_{i=1}^{3} w_i x_i$. (This last kind of sum, a \emph{weighted sum},
is what a single artificial neuron computes.)
\end{solution}

\begin{exercise}
Let $x = (1, 4, 2, 5)$ and $y = (3, 0, 1, 2)$, meaning $x_1 = 1$, $y_1 = 3$,
and so on. Compute $\sum_i x_i$, $\sum_i y_i$, $\sum_i x_i y_i$ and
$\left(\sum_i x_i\right)\left(\sum_i y_i\right)$.
\end{exercise}
\begin{solution}
$\sum_i x_i = 12$, $\sum_i y_i = 6$,
$\sum_i x_i y_i = 3 + 0 + 2 + 10 = 15$, and
$\left(\sum_i x_i\right)\left(\sum_i y_i\right) = 72$. The last two differ,
confirming that a sum of products is not a product of sums.
\end{solution}

\begin{exercise}
Use the rules for sums to show that the deviations from the mean always add up
to zero: $\sum_{i=1}^n (x_i - \bar x) = 0$.
\end{exercise}
\begin{solution}
Split the sum and use $\sum_{i=1}^n c = nc$ with the constant $c = \bar x$:
\[
  \sum_{i=1}^n (x_i - \bar x) = \sum_{i=1}^n x_i - n\bar x
  = n\bar x - n\bar x = 0,
\]
because $\sum_{i=1}^n x_i = n\bar x$ by the definition of the mean.
\end{solution}

\begin{exercise}
True values are $y = (2, 4, 6)$ and a model predicts $\hat y = (3, 4, 4)$.
Compute the mean squared error and the mean absolute error
$\tfrac1n\sum_i \abs{y_i - \hat y_i}$. (The absolute value $\abs{\cdot}$ simply
removes a minus sign; see \cref{sec:absolute-value}.)
\end{exercise}
\begin{solution}
Errors $y_i - \hat y_i$: $-1, 0, 2$. Squares: $1, 0, 4$, so
$\MSE = \tfrac53 \approx 1.67$. Absolute errors: $1, 0, 2$, so the mean
absolute error is $\tfrac33 = 1$.
\end{solution}

%----------------------------------------------------------------------
\section{Solving equations}\label{sec:equations}
%----------------------------------------------------------------------

An \term{equation} states that two expressions are equal, for example
$3x + 5 = 20$. A \term[solution of an equation]{solution} is a value of the
variable that makes the statement true. Here $x = 5$ is a solution, because
$3 \times 5 + 5 = 20$.

\subsection{The balance principle}

\begin{picturethis}
Imagine an old-fashioned balance scale with $3x + 5$ in the left pan and $20$
in the right pan, perfectly level. If you remove $5$ from \emph{both} pans, it
stays level: $3x = 15$. If you then divide the contents of \emph{both} pans
into three equal parts and keep one part of each, it still stays level:
$x = 5$. Solving an equation means doing the same thing to both sides until
the variable stands alone.
\end{picturethis}

The allowed moves are: add or subtract the same number on both sides; multiply
or divide both sides by the same \emph{non-zero} number; and simplify either
side. Each move undoes one operation that was applied to $x$, in reverse order.

\begin{example}[Linear equations]\label{ex:linear-equations}
Solve (a) $3x + 5 = 20$; (b) $5x - 3 = 2x + 9$; (c) $\tfrac{x}{4} + 1 = 3$;
(d) $\tfrac{x+1}{3} = \tfrac{x-1}{2}$.

\sol
(a) Subtract $5$: $3x = 15$. Divide by $3$: $x = 5$.
\emph{Check:} $3\times5 + 5 = 20$. \checkmark

(b) Gather the $x$ terms on one side by subtracting $2x$: $3x - 3 = 9$. Add
$3$: $3x = 12$. Divide by $3$: $x = 4$.
\emph{Check:} left side $5\times4 - 3 = 17$, right side $2\times4+9 = 17$.
\checkmark

(c) Subtract $1$: $\tfrac{x}{4} = 2$. Multiply by $4$: $x = 8$.

(d) Clear the fractions by multiplying both sides by $6$, a common multiple of
the denominators: $2(x+1) = 3(x-1)$. Expand: $2x + 2 = 3x - 3$. Subtract $2x$
and add $3$: $5 = x$. \emph{Check:} $\tfrac{6}{3} = 2 = \tfrac42$. \checkmark
\end{example}

Always check your answer by substituting it back into the \emph{original}
equation. It takes seconds and catches most mistakes.

Equations of the form $ax + b = c$ are called \term[linear equation]{linear
equations}, because (as we will see in \cref{ch:functions}) their graphs are
straight lines. Every linear equation with $a \neq 0$ has exactly one
solution, $x = \tfrac{c-b}{a}$. Some equations that look linear have no
solution at all ($x + 1 = x + 2$ simplifies to the false statement $1 = 2$) or
are true for every $x$ ($2(x+1) = 2x + 2$ simplifies to $2 = 2$; such an
equation is called an \emph{identity}).

\subsection{Rearranging formulas}

The same moves let us solve a formula for any one of its variables. To convert
Fahrenheit to Celsius, solve $F = \tfrac95 C + 32$ for $C$: subtract $32$, then
multiply by $\tfrac59$:
\[
  C = \tfrac59(F - 32).
\]
A body temperature of $98.6\,^\circ$F is therefore
$\tfrac59 \times 66.6 = 37\,^\circ$C.

\subsection{Two equations in two unknowns}

Sometimes we have two unknowns and two facts about them. For example: two
numbers add up to $10$ and differ by $2$. What are they? Writing $x$ and $y$
for the numbers,
\[
  \begin{cases}
    x + y = 10,\\
    x - y = 2.
  \end{cases}
\]
This is a \term[system of linear equations]{system of linear equations}. Two
standard methods solve it.

\textbf{Elimination.} Add the two equations: the $y$'s cancel and we get
$2x = 12$, so $x = 6$. Substituting into the first equation gives $y = 4$.

\textbf{Substitution.} Solve one equation for one variable and substitute
into the other. From the second equation, $x = y + 2$; substituting into the
first gives $(y + 2) + y = 10$, so $2y = 8$, $y = 4$, and $x = 6$.

\begin{example}[Fitting a line through two points]\label{ex:line-two-points}
Find numbers $w$ and $b$ so that the formula $y = wx + b$ gives $y = 3$ when
$x = 1$ and gives $y = 7$ when $x = 3$.

\sol Substituting the two data points gives two equations:
\[
  \begin{cases}
    w + b = 3,\\
    3w + b = 7.
  \end{cases}
\]
Subtracting the first equation from the second eliminates $b$: $2w = 4$, so
$w = 2$. Then $b = 3 - w = 1$. The formula is $y = 2x + 1$.
\emph{Check:} $2\times1+1 = 3$ and $2\times3+1 = 7$. \checkmark
\end{example}

\begin{mlnote}[From two equations to millions]
\Cref{ex:line-two-points} is a miniature version of training a model: we
chose the parameters $w$ and $b$ so that the model reproduces the data. With
two data points and two parameters we can fit the data exactly. Real data sets
have many more data points than parameters, so there are more equations than
unknowns and usually no exact solution. We will then look for the parameters
that make the errors as small as possible, a problem called \emph{least
squares} (\cref{ch:linalg2}). Large systems of equations are written and
solved with matrices (\cref{ch:linalg1}).
\end{mlnote}

\sectionexercises

\begin{exercise}
Solve and check: (a) $4x - 7 = 13$; (b) $2(x+3) = 5x - 3$;
(c) $\tfrac{x}{3} - 2 = 4$; (d) $\tfrac{2x-1}{5} = \tfrac{x+2}{3}$.
\end{exercise}
\begin{solution}
(a) $4x = 20$, $x = 5$.
(b) $2x + 6 = 5x - 3$, so $9 = 3x$ and $x = 3$; check: $12 = 12$.
(c) $\tfrac{x}{3} = 6$, $x = 18$.
(d) Multiply by $15$: $3(2x - 1) = 5(x + 2)$, so $6x - 3 = 5x + 10$ and
$x = 13$; check: $\tfrac{25}{5} = 5 = \tfrac{15}{3}$.
\end{solution}

\begin{exercise}
Rearrange each formula for the indicated variable: (a) $v = u + at$ for $t$;
(b) $A = \tfrac12 bh$ for $h$; (c) $z = \tfrac{x - \mu}{\sigma}$ for $x$.
(The letters $\mu$ and $\sigma$ are the Greek letters ``mu'' and ``sigma''.
The quantity $z$ is the \emph{z-score} of statistics, which you will meet in
\cref{ch:probability}.)
\end{exercise}
\begin{solution}
(a) $t = \tfrac{v - u}{a}$ (assuming $a \neq 0$).
(b) $h = \tfrac{2A}{b}$ (assuming $b \neq 0$).
(c) Multiply by $\sigma$: $z\sigma = x - \mu$, so $x = \mu + \sigma z$.
\end{solution}

\begin{exercise}
Solve the systems: (a) $x + y = 7$, $x - y = 3$; (b) $3x + 2y = 16$,
$x - 2y = 0$; (c) $2x + y = 5$, $4x + 2y = 10$; (d) $x + y = 1$,
$x + y = 3$.
\end{exercise}
\begin{solution}
(a) Adding: $2x = 10$, so $x = 5$, $y = 2$.
(b) Adding: $4x = 16$, so $x = 4$ and then $y = 2$.
(c) The second equation is twice the first, so they carry the same
information: every pair with $y = 5 - 2x$ is a solution (infinitely many
solutions).
(d) The two equations contradict each other ($1 \neq 3$), so there is no
solution.
\end{solution}

\begin{exercise}
Find the line $y = wx + b$ that passes through the points $(x,y) = (2,5)$ and
$(4,9)$.
\end{exercise}
\begin{solution}
$2w + b = 5$ and $4w + b = 9$. Subtracting: $2w = 4$, so $w = 2$ and
$b = 5 - 4 = 1$. The line is $y = 2x + 1$.
\end{solution}

\begin{exercise}
A theatre sells adult tickets at $12$ and child tickets at $7$. A group buys
$20$ tickets for a total of $190$. How many adult tickets did they buy?
\end{exercise}
\begin{solution}
Let $a$ and $c$ be the numbers of adult and child tickets: $a + c = 20$ and
$12a + 7c = 190$. Substituting $c = 20 - a$: $12a + 140 - 7a = 190$, so
$5a = 50$ and $a = 10$ (and $c = 10$).
\end{solution}

%----------------------------------------------------------------------
\section{Inequalities}\label{sec:inequalities}
%----------------------------------------------------------------------

\subsection{Inequality symbols}

An \term{inequality} compares two quantities that need not be equal:
\begin{center}
\renewcommand{\arraystretch}{1.2}
\begin{tabular}{@{}lll@{}}
\toprule
\textbf{Symbol} & \textbf{Read as} & \textbf{Example}\\
\midrule
$a < b$ & $a$ is less than $b$ & $-3 < 2$\\
$a > b$ & $a$ is greater than $b$ & $7 > 4$\\
$a \le b$ & $a$ is less than or equal to $b$ & $5 \le 5$\\
$a \ge b$ & $a$ is greater than or equal to $b$ & $0.3 \ge 0$\\
$a \neq b$ & $a$ is not equal to $b$ & $1 \neq 2$\\
\bottomrule
\end{tabular}
\end{center}
On the number line, $a < b$ means that $a$ lies to the left of $b$. (A memory
aid: the symbol always opens towards the larger number.)

\subsection{Solving inequalities}

Inequalities are solved almost exactly like equations: you may add or
subtract the same number on both sides, and multiply or divide both sides by
the same \emph{positive} number. There is one crucial exception.

\begin{theorem}[Multiplying an inequality by a negative number]
\label{thm:flip-inequality}
If $a < b$ and $c < 0$, then $ca > cb$. In words: multiplying or dividing both
sides of an inequality by a negative number \emph{reverses} the direction of
the inequality.
\end{theorem}

\begin{intuition}
Multiplying by $-1$ reflects the number line in the point $0$, like a mirror.
Since $2 < 5$, the point $2$ is to the left of $5$; after reflection, $-2$ is to
the \emph{right} of $-5$, so $-2 > -5$. Order is reversed by the mirror.
\end{intuition}

\begin{example}[Solving an inequality]\label{ex:inequality}
Solve $3 - 2x \ge 7$.

\sol Subtract $3$ from both sides: $-2x \ge 4$. Divide by $-2$, which is
negative, so the inequality flips: $x \le -2$. The solution set is the
interval $(-\infty, -2]$:
\begin{center}
\begin{tikzpicture}[x=0.6cm]
  \draw[{Stealth[length=2mm]}-{Stealth[length=2mm]},thin,gray] (-6.5,0) -- (3.5,0);
  \foreach \x in {-6,...,3}{
    \draw[gray] (\x,0.07) -- (\x,-0.07) node[below,font=\scriptsize] {$\x$};}
  \draw[line width=2.5pt,PlotBlue,{Stealth[length=3mm]}-] (-6.5,0) -- (-2,0);
  \fill[PlotBlue] (-2,0) circle (3.5pt);
\end{tikzpicture}
\end{center}
\emph{Check} a value inside the set, say $x = -3$: $3 + 6 = 9 \ge 7$
\checkmark; and one outside, say $x = 0$: $3 \ge 7$ is false. \checkmark
\end{example}

A \emph{compound inequality} such as $-1 < 2x + 3 \le 9$ is solved by
operating on all three parts at once: subtract $3$ to get $-4 < 2x \le 6$, then
divide by $2$ to get $-2 < x \le 3$, i.e.\ the interval $(-2, 3]$.

\subsection{Maximum and minimum}

The \term{maximum} of two numbers is the larger of them, and the
\term{minimum} is the smaller:
\[
  \max(3, 7) = 7, \qquad \min(3, 7) = 3, \qquad \max(-2, -5) = -2 .
\]
More generally, $\max$ and $\min$ of a list of numbers pick out its largest
and smallest elements.

\begin{mlnote}[Inequalities in machine learning]
Inequalities express constraints: a probability $p$ always satisfies
$0 \le p \le 1$, and a learning rate must be positive. The maximum function
also defines the most popular building block of modern neural networks, the
\term{rectified linear unit}, or ReLU:
\[
  \ReLU(x) = \max(0, x) .
\]
It passes positive numbers through unchanged and replaces negative numbers by
$0$. So $\ReLU(2.5) = 2.5$ and $\ReLU(-1.3) = 0$. Another example is
\emph{clipping}, which forces a number into an interval $[-c, c]$ using
$\min(\max(x, -c), c)$; ``gradient clipping'' uses this idea to stop training
from blowing up.
\end{mlnote}

\sectionexercises

\begin{exercise}
Solve, and write the answer in interval notation: (a) $2x + 5 < 11$;
(b) $4 - 3x \le 10$; (c) $-3 \le 2x - 1 < 7$.
\end{exercise}
\begin{solution}
(a) $2x < 6$, so $x < 3$: $(-\infty, 3)$.
(b) $-3x \le 6$; dividing by $-3$ flips the sign: $x \ge -2$:
$[-2, \infty)$.
(c) Add $1$: $-2 \le 2x < 8$; divide by $2$: $-1 \le x < 4$: $[-1, 4)$.
\end{solution}

\begin{exercise}
Evaluate: (a) $\max(-2, 5)$; (b) $\min(-2, 5)$; (c) $\ReLU(-3.7)$;
(d) $\ReLU(2.1)$; (e) $\max(1, 4, -6, 2)$.
\end{exercise}
\begin{solution}
(a) $5$; (b) $-2$; (c) $0$; (d) $2.1$; (e) $4$.
\end{solution}

\begin{exercise}
For which values of $x$ is $\ReLU(x) = x$? For which is $\ReLU(x) = 0$?
\end{exercise}
\begin{solution}
$\ReLU(x) = x$ exactly when $x \ge 0$; $\ReLU(x) = 0$ exactly when $x \le 0$
(both hold at $x = 0$).
\end{solution}

\begin{exercise}
The clipping function with $c = 1$ is $\operatorname{clip}(x) =
\min(\max(x, -1), 1)$. Compute $\operatorname{clip}(-3)$,
$\operatorname{clip}(0.5)$ and $\operatorname{clip}(2)$, and describe in words
what the function does.
\end{exercise}
\begin{solution}
$\operatorname{clip}(-3) = \min(-1, 1) = -1$;
$\operatorname{clip}(0.5) = \min(0.5, 1) = 0.5$;
$\operatorname{clip}(2) = \min(2, 1) = 1$. It leaves numbers in $[-1,1]$
unchanged and replaces anything outside by the nearest endpoint.
\end{solution}

\begin{exercise}
Starting from the true statement $3 < 8$, multiply both sides by $-2$. Is
$-6 < -16$ true? Explain what goes wrong and how to fix it.
\end{exercise}
\begin{solution}
$-6 < -16$ is false: $-6$ lies to the right of $-16$ on the number line.
Multiplying by a negative number reverses the inequality, so the correct
statement is $-6 > -16$.
\end{solution}

%----------------------------------------------------------------------
\section{Absolute value}\label{sec:absolute-value}
%----------------------------------------------------------------------

\begin{definition}[Absolute value]\label{def:absolute-value}
The \term{absolute value} of a real number $x$ is
\[
  \abs{x} = \begin{cases}
    x & \text{if } x \ge 0,\\
    -x & \text{if } x < 0.
  \end{cases}
\]
\end{definition}

In practice: the absolute value removes a minus sign, if there is one.
$\abs{5} = 5$, $\abs{-5} = 5$, $\abs{0} = 0$. (In the second case the
definition says $\abs{-5} = -(-5) = 5$.)

\begin{intuition}[Absolute value is distance]
$\abs{x}$ is the distance from $x$ to $0$ on the number line, and distances are
never negative. More usefully, $\abs{a - b}$ is the distance between the
numbers $a$ and $b$ (\cref{fig:abs-distance}). For example, the distance
between $-2$ and $3$ is $\abs{-2 - 3} = \abs{-5} = 5$, and it does not matter
in which order we subtract: $\abs{3 - (-2)} = 5$ too.
\end{intuition}

\begin{figure}[htbp]
\centering
\begin{tikzpicture}[x=1cm]
  \draw[{Stealth[length=2mm]}-{Stealth[length=2mm]},thick] (-4.5,0) -- (4.5,0);
  \foreach \x in {-4,...,4}
    \draw (\x,0.07) -- (\x,-0.07) node[below,font=\small] {$\x$};
  \fill[PlotRed] (-2,0) circle (2.5pt) node[above=4pt] {$a=-2$};
  \fill[PlotBlue] (3,0) circle (2.5pt) node[above=4pt] {$b=3$};
  \draw[decorate,decoration={brace,amplitude=6pt,mirror},thick]
    (-2,-0.55) -- (3,-0.55) node[midway,below=8pt] {$\abs{a-b} = 5$};
\end{tikzpicture}
\caption{The distance between two numbers $a$ and $b$ on the number line is
  $\abs{a-b}$.}
\label{fig:abs-distance}
\end{figure}

\begin{proposition}[Properties of absolute value]\label{prop:abs-properties}
For all real numbers $a$ and $b$:
\begin{enumerate}
  \item $\abs{a} \ge 0$, and $\abs{a} = 0$ only when $a = 0$;
  \item $\abs{ab} = \abs{a}\,\abs{b}$ and, if $b \neq 0$,
    $\abs{a/b} = \abs{a}/\abs{b}$;
  \item $\abs{a + b} \le \abs{a} + \abs{b}$ (the \term{triangle inequality});
  \item $\sqrt{a^2} = \abs{a}$.
\end{enumerate}
\end{proposition}

The triangle inequality says that going from $0$ to $a + b$ can never be
longer than going first a distance $\abs{a}$ and then a distance $\abs{b}$.
Equality holds when $a$ and $b$ have the same sign; when the signs differ, the
two steps partly cancel: $\abs{3 + (-5)} = 2$, while $\abs3 + \abs{-5} = 8$.
The name will make more sense in \cref{ch:linalg1}, where the same inequality
holds for the sides of a triangle.

\subsection{Equations and inequalities with absolute values}

Reading absolute values as distances makes these problems easy.

\begin{example}[Absolute value equations and inequalities]
\label{ex:abs-equations}
Solve (a) $\abs{x - 3} = 2$; (b) $\abs{x - 3} < 2$; (c) $\abs{x-3} \ge 2$;
(d) $\abs{2x + 1} \le 5$.

\sol (a) ``The distance from $x$ to $3$ is $2$'': $x = 1$ or $x = 5$.

(b) ``The distance from $x$ to $3$ is less than $2$'': $1 < x < 5$, the
interval $(1, 5)$.

(c) ``At least $2$ away from $3$'': $x \le 1$ or $x \ge 5$.

(d) $\abs{u} \le 5$ means $-5 \le u \le 5$. With $u = 2x + 1$:
$-5 \le 2x + 1 \le 5$, so $-6 \le 2x \le 4$, i.e.\ $-3 \le x \le 2$.
\end{example}

\begin{pitfall}[Absolute value does not split over sums]
$\abs{a + b}$ is not in general equal to $\abs a + \abs b$ (see the example
above). Also, $-\abs{x}$ is always $\le 0$, while $\abs{-x} = \abs{x}$ is
always $\ge 0$; these are different things.
\end{pitfall}

\begin{mlnote}[Two ways of measuring errors]
If a model predicts $\hat y$ and the true value is $y$, the size of the error is
$\abs{y - \hat y}$, regardless of whether the prediction was too high or too
low. Averaging over a data set gives the \term{mean absolute error},
\[
  \text{MAE} = \frac1n \sum_{i=1}^{n} \abs{y_i - \hat y_i},
\]
an alternative to the mean squared error of \cref{sec:summation}. Squaring
punishes large errors much more than small ones (an error of $10$ contributes
$100$ to the MSE but only $10$ to the MAE), so the two measures lead to models
with different behaviour. In \cref{ch:linalg2} absolute values reappear in the
``$L^1$ norm'', which is used to encourage models to set many weights to
exactly zero.
\end{mlnote}

\sectionexercises

\begin{exercise}
Evaluate: $\abs{-8}$, $\abs{3 - 10}$, $\abs{-2}\cdot\abs{5}$, $\abs{-4 + 9}$,
$-\abs{-6}$.
\end{exercise}
\begin{solution}
$8$; $\abs{-7} = 7$; $2\cdot5 = 10$; $\abs{5} = 5$; $-6$.
\end{solution}

\begin{exercise}
Solve: (a) $\abs{x+2} = 5$; (b) $\abs{x-1} \le 3$; (c) $\abs{2x - 4} > 6$.
\end{exercise}
\begin{solution}
(a) $x + 2 = 5$ or $x + 2 = -5$, so $x = 3$ or $x = -7$.
(b) $-3 \le x - 1 \le 3$, so $-2 \le x \le 4$: $[-2, 4]$.
(c) $2x - 4 > 6$ or $2x - 4 < -6$, so $x > 5$ or $x < -1$.
\end{solution}

\begin{exercise}
What is the distance between $-3.5$ and $2$ on the number line?
\end{exercise}
\begin{solution}
$\abs{-3.5 - 2} = \abs{-5.5} = 5.5$.
\end{solution}

\begin{exercise}
Check the triangle inequality for $a = 4$, $b = -7$. For which pairs of
numbers does $\abs{a+b} = \abs a + \abs b$ hold exactly?
\end{exercise}
\begin{solution}
$\abs{4 + (-7)} = 3 \le \abs4 + \abs{-7} = 11$. Equality holds exactly when
$a$ and $b$ do not have opposite signs, i.e.\ when both are $\ge 0$ or both are
$\le 0$.
\end{solution}

\begin{exercise}
True values are $y = (10, 12, 8)$ and predictions are $\hat y = (11, 9, 8)$.
Compute the MAE and the MSE. Which data point dominates the MSE, and why?
\end{exercise}
\begin{solution}
Errors: $-1, 3, 0$. MAE $= \tfrac{1 + 3 + 0}{3} = \tfrac43 \approx 1.33$.
MSE $= \tfrac{1 + 9 + 0}{3} = \tfrac{10}{3} \approx 3.33$. The second data
point contributes $9$ of the $10$ units of squared error: squaring makes the
largest error dominate.
\end{solution}

%----------------------------------------------------------------------
\section{Chapter summary}\label{sec:ch1-summary}
%----------------------------------------------------------------------

\begin{keytakeaways}
\begin{itemize}
  \item The number systems are nested: $\N \subset \Z \subset \Q \subset \R$.
    Machine learning works almost entirely with real numbers, which are
    stored approximately by computers.
  \item The commutative, associative and distributive laws let us rearrange
    calculations. Never divide by zero.
  \item Fractions: find a common denominator to add; multiply straight across;
    divide by multiplying by the reciprocal. Always ask what a percentage is a
    percentage \emph{of}.
  \item Exponent laws: $a^m a^n = a^{m+n}$, $(a^m)^n = a^{mn}$, $a^0 = 1$,
    $a^{-n} = 1/a^n$, $a^{1/n} = \sqrt[n]{a}$. But $(a+b)^2 \neq a^2 + b^2$.
  \item Order of operations: brackets, exponents, multiplication and division
    (left to right), addition and subtraction (left to right).
  \item Variables are named boxes for numbers; a model is an expression with
    inputs and parameters (weights and biases).
  \item $\sum_{i=1}^n x_i$ means $x_1 + \dots + x_n$. Constants can be pulled
    out of sums and sums can be split, but a sum of products is not a product
    of sums. The mean is $\bar x = \frac1n\sum_i x_i$.
  \item Solve equations by doing the same thing to both sides. Two unknowns
    need two equations (elimination or substitution).
  \item Inequalities work like equations, except that multiplying by a
    negative number reverses them. $\ReLU(x) = \max(0, x)$.
  \item $\abs{a - b}$ is the distance between $a$ and $b$; mean absolute error
    and mean squared error are two ways to measure a model's mistakes.
\end{itemize}
\end{keytakeaways}

In the next chapter we turn expressions into \emph{functions}: machines that
take an input and produce an output. Functions are the mathematical name for
what a machine learning model is.
